% ===================================================================
%  اللوحة رقم ٦ — البيت من الداخل. — الأثاث. — الطعام. — الإضاءة.
%  Traduction 2026 d'apres texto/io, controlee sur texto/fr, texto/en
%  et texto/es. Consigne : voir texto/ar/10-lawha-01.tex.
%
%  Deux notes, toutes deux marquees « (*) », et deux appels : celui de
%  « babo » au bloc c1-12-1, celui de « lambrequino » au bloc c2-01-1.
%  L'ordre les departage, comme dans les autres colonnes.
%
%  LA FEMME DE CHAMBRE. L'ido du bloc c1-06-1 porte « (6) » la ou le
%  francais porte « (16) », et c'est le francais qui a raison. L'arabe
%  suit le francais, comme les autres traductions.
% ===================================================================

%%K t06-apar-1 apar
\VUsaut{6.0mm}
\VUtitre{15.0pt}{60}{اللوحة رقم ٦}
\VUsaut{1.4mm}
\VUtitre{11.4pt}{0}{البيت من الداخل، الأثاث، الطعام،}
\VUsaut{0.4mm}
\VUtitre{11.4pt}{0}{الإضاءة.}
\VUsaut{2.2mm}

%%K t06-apar-2 apar
تمّ البيت. وانتقل عمّ يوحنا إلى مسكنه الجديد الذي نعرف تقسيمه من قبلُ.
فلنرَ الآن كيف أثّثه. سنزور أولًا:

%%K t06-tit-1 sub
\VUpk{10.2pt}{40}{غرفة النوم.}
\VUsaut{3.0mm}

%%K t06-c1-01-1 p
1. --- جئنا في وقت غير مناسب، فالخادمة مشغولة بتنظيف \VUgras{الغرفة}.

%%K t06-c1-02-1 p
2. --- على \VUgras{المغسلة} \textsuperscript{(1)} تتناثر هنا وهناك الأشياء التي
يُغتسل بها ويُمشَّط ويُتزيَّن. وهي \VUgras{الطست} \textsuperscript{(2)}
و\VUgras{إبريق الماء} \textsuperscript{(3)}، و\VUgras{فرشاة الشعر}
\textsuperscript{(4)}، و\VUgras{المشط} \textsuperscript{(5)}، و\VUgras{الصابون}
\textsuperscript{(6)}، و\VUgras{الإسفنجة} \textsuperscript{(7)}، وأخيرًا
\VUgras{قارورة} \textsuperscript{(8)} صغيرة لغسول الفم.

%%K t06-c1-03-1 p
3. --- وعلى الأرض، أمام المغسلة، ها هو \VUgras{الدلو} \textsuperscript{(9)} للماء
الوسخ.

%%K t06-c1-04-1 p
4. --- وفي \VUgras{الغرفة الخارجية} \textsuperscript{(10)} نرى بضعة
\VUgras{مشاجب} \textsuperscript{(13)} و\VUgras{مظلتين} \textsuperscript{(11)} في
\VUgras{حامل المظلات} \textsuperscript{(12)}.

%%K t06-c1-05-1 p
5. --- وعلى الجانب الآخر من الباب تقوم \VUgras{الخزانة ذات المرآة}
\textsuperscript{(14)}، وفيها \VUgras{البياضات} \textsuperscript{(15)}.

%%K t06-c1-06-1 p
6. --- ولنغتنم انشغال \VUgras{وصيفة الغرف} \textsuperscript{(16)} بترتيب
\VUgras{السرير} \textsuperscript{(17)} لننظر في أجزائه. \VUgras{هيكل السرير}
\textsuperscript{(20)} يحمل \VUgras{فراشًا زنبركيًا} \textsuperscript{(21)} جيدًا،
ستضع عليه جوستين بعد قليل \VUgras{مرتبة} \textsuperscript{(22)} صوفية. ثم تأخذ
عن الكراسي \VUgras{الملاءتين} \textsuperscript{(23)} و\VUgras{البطانية}
\textsuperscript{(24)}. ثم تضع عند رأس السرير \VUgras{المسند} \textsuperscript{(25)}
و\VUgras{الوسادة} \textsuperscript{(26)}، وهما موضوعان مع \VUgras{اللحاف}
\textsuperscript{(27)} على \VUgras{سجادة السرير} \textsuperscript{(32)}. وتنفض
\VUgras{الستائر} \textsuperscript{(19)} و\VUgras{الناموسية} \textsuperscript{(18)}
بمنفضة الريش، فيتم شطر من العمل.

%%K t06-c1-07-1 p
7. --- ثم يبقى عليها أن ترتب \VUgras{المنضدة الجانبية} \textsuperscript{(28)}
قليلًا، وتملأ زنبرك \VUgras{المنبّه} \textsuperscript{(29)}، وتهيّئ
\VUgras{مصباح الليل} \textsuperscript{(30)}، وترفع \VUgras{الشمعة}
\textsuperscript{(31)} لتنظف الشمعدان.

%%K t06-tit-2 sub
\VUsaut{5.0mm}
\VUpk{10.2pt}{40}{سلة الخياطة وأدوات المدفأة.}
\VUsaut{3.0mm}

%%K t06-c1-08-1 p
8. --- و\VUgras{سلة الخياطة} \textsuperscript{(34)} إلى جانب \VUgras{المقعد}
\textsuperscript{(33)} في حاجة شديدة إلى الترتيب هي أيضًا. سيلزمها أن تطوي
\VUgras{التطريز} \textsuperscript{(35)}، وأن تضع في \VUgras{منضدة الخياطة}
\textsuperscript{(38)} \VUgras{الكشتبان} و\VUgras{الخيط} \textsuperscript{(36)}
و\VUgras{الإبر} \textsuperscript{(37)} و\VUgras{المقص} \textsuperscript{(39)}، وأن
تقفل غطاء المنضدة بـ\VUgras{المفتاح} \textsuperscript{(40)}.

%%K t06-c1-09-1 p
9. --- وعلى \VUgras{المنضدة المستديرة} \textsuperscript{(41)} في الوسط، ستبدّل
جوستين ماء \VUgras{الدورق} \textsuperscript{(42)}. وستضع \VUgras{السكر} في
\VUgras{السكرية} \textsuperscript{(43)}، وتنظف \VUgras{ملقط السكر}
\textsuperscript{(44)}، وترفع \VUgras{فرشاة الثياب} \textsuperscript{(46)}.

%%K t06-c1-10-1 p
10. --- والجانب الأيمن من الغرفة سيكلفها عملًا أقل، فـ\VUgras{الأريكة}
\textsuperscript{(47)} و\VUgras{وسادتها} \textsuperscript{(48)} في موضعهما، وما دام
الجو غير بارد فلم يُحرَّك شيء من أدوات \VUgras{المدفأة} \textsuperscript{(49)}.
\VUgras{المجرفة} \textsuperscript{(50)} و\VUgras{الملقط} \textsuperscript{(51)}
و\VUgras{حاملا الحطب} \textsuperscript{(55)} و\VUgras{حاجز الرماد}
\textsuperscript{(56)} نظيفة؛ و\VUgras{حاجب النار} \textsuperscript{(54)} لم يُنقَل،
و\VUgras{صندوق الحطب} \textsuperscript{(52)} ملآن بـ\VUgras{الحطب}
\textsuperscript{(53)}.

%%K t06-noto-1 noto
\VUnotes{143.2mm}{%
(*) Babo = طفل صغير؛ بالإنجليزية \textit{baby}، وبالفرنسية
\textit{bébé}، وبالإيطالية \textit{bambino}.%
}

%%K t06-noto-2 noto
\VUnotes{143.2mm}{%
(*) Lambrequino = بالفرنسية \textit{lambrequin}، وبالإسبانية
\textit{lambrequines}، وبالبرتغالية \textit{lambrequins}، بل
وبالألمانية والإنجليزية \textit{lambrequins} --- فهي واحدة في الألمانية
والإنجليزية والفرنسية والبرتغالية والإسبانية.%
}

%%K t06-c1-11-1 p
11. --- ويبقى عليها أن تنفض \VUgras{ساعة الرف} \textsuperscript{(57)}
و\VUgras{الشمعدانات} \textsuperscript{(58)} و\VUgras{صواني الحلي}
\textsuperscript{(59)}، ثم تمسح \VUgras{المرآة} \textsuperscript{(60)}.

%%K t06-c1-12-1 p
12. --- أما \VUgras{مهد} \textsuperscript{(61)} الرضيع (الـ babo (*)) فهو جاهز
من قبلُ.

%%K t06-tit-3 sub
\VUsaut{5.0mm}
\VUpk{10.2pt}{40}{صالة الاستقبال.}
\VUsaut{3.0mm}

%%K t06-c2-01-1 p
1. --- وها هي الآن صالة الاستقبال. وهي \VUgras{غرفة} جميلة جدًا. والمدفأة
في الركن الأيمن مزينة بـ\VUgras{تمثال صغير} \textsuperscript{(1)} رقيق
و\VUgras{مزهريتين} \textsuperscript{(3)} جميلتين، ومسدل عليها \VUgras{شراعة}
\textsuperscript{(2)} من المخمل (*).

%%K t06-c2-02-1 p
2. --- ولئلا يشعل شرر الموقد السجادة، وُضع أمامها \VUgras{حاجز نحاسي}
\textsuperscript{(4)}.

%%K t06-c2-03-1 p
3. --- وإلى جانبها ينتصب \VUgras{مكتب} \textsuperscript{(5)} أنيق للسيدات وهو
مفتوح. وفيه تكتب \VUgras{ربّة البيت} \textsuperscript{(23)} رسائلها.

%%K t06-c2-04-1 p
4. --- والنافذة مزيّنة بـ\VUgras{ستائر التُّل} \textsuperscript{(6)} تشدّها
\VUgras{شرائط} \textsuperscript{(8)} معلقة في \VUgras{الخطاطيف}
\textsuperscript{(7)}، وبستائر قماشية ثقيلة يعلوها \VUgras{كرنيش}
\textsuperscript{(9)}.

%%K t06-c2-05-1 p
5. --- وفي تجويف النافذة تحمل \VUgras{مزهرية أرضية} \textsuperscript{(11)}
\VUgras{نبتة} \textsuperscript{(10)} خضراء، نخلة رائعة تكاد أوراقها تبلغ
\VUgras{مصباح الزيت} \textsuperscript{(12)}.

%%K t06-c2-06-1 p
6. --- و\VUgras{غطاء المصباح} \textsuperscript{(13)} صنعته ماري،
\VUgras{الآنسة} \textsuperscript{(14)} الفاتنة الجالسة إلى \VUgras{البيانو}
\textsuperscript{(15)}. وهي، جالسة معتدلة على \VUgras{المقعد} \textsuperscript{(16)}،
تتمرن على قطعة. وهي ماهرة جدًا ومرتبة جدًا، وشاهد ذلك \VUgras{خزانة
النوتات} \textsuperscript{(17)}: فهي في نظام تام.

%%K t06-tit-4 sub
\VUsaut{5.0mm}
\VUpk{10.2pt}{40}{الشقيّ.}
\VUsaut{3.0mm}

%%K t06-c2-07-1 p
7. --- وأخوها بولس يبحث عن \VUgras{كتاب} \textsuperscript{(19)} في
\VUgras{خزانة الكتب} \textsuperscript{(18)}. وهو \VUgras{ولد} \textsuperscript{(20)}
كثير الحركة لا يُطاق. وإذ يتعلق بالخزانة ليبلغ الكتاب الموضوع في علوّ
شديد، يوشك أن يُسقط \VUgras{التمثال النصفي} \textsuperscript{(21)} الذي فوقها.

%%K t06-c2-08-1 p
8. --- وهو يغتنم لهذه الحماقة انشغال أمه بالحديث مع صديقة، \VUgras{سيدة}
\textsuperscript{(22)} جاءت لتقضي عندها أيامًا. والأم جالسة على \VUgras{الكنبة}
\textsuperscript{(24)} إلى جانب \VUgras{المقعد الوثير} \textsuperscript{(25)}،
وصديقتها على \VUgras{الكرسي} \textsuperscript{(27)} الذي لا نرى منه إلا
\VUgras{ظهره} \textsuperscript{(26)}.

%%K t06-c2-09-1 p
9. --- و\VUgras{الإطار} \textsuperscript{(30)} الكبير المعلق على الجدار يحوي
\VUgras{صورة} \textsuperscript{(29)} أحد الأقارب. وفي \VUgras{الألبوم}
\textsuperscript{(32)} الموضوع على المنضدة جُمعت \VUgras{صور} \textsuperscript{(33)}
بقية الأسرة. وقد كان الأولاد يقلّبونه منذ قليل، فـ\VUgras{الكراسي}
\textsuperscript{(31)} لا تزال مجتمعة حول المنضدة.

%%K t06-c2-10-1 p
10. --- و\VUgras{المزهرية الصينية} \textsuperscript{(34)} الخزفية قرب
\VUgras{سكين الورق} \textsuperscript{(35)} أُهديت إلى ماري في عيد اسمها، أما
\VUgras{الحاجز الطيّاء} \textsuperscript{(36)} الذي يملأ ركن الغرفة فقد جلبه
عمها من الصين.

%%K t06-c2-11-1 p
11. --- وهذه الصالة تبدو لطيفة جدًا، تزينها \VUgras{اللوحات}
\textsuperscript{(37)} الجميلة المعلقة على \VUgras{الجدار} \textsuperscript{(42)}،
وتضيئها \VUgras{الثريا} \textsuperscript{(38)} المعلقة في السقف، بـ\VUgras{مصابيحها
الكهربائية} \textsuperscript{(39)} التي تشعّ في المساء إشعاعًا مبهجًا.
و\VUgras{السجادة} \textsuperscript{(40)} الوثيرة المفروشة على الأرض،
و\VUgras{الجرس الكهربائي} \textsuperscript{(41)} المريح جدًا، يجعلانها في غاية
الراحة.

%%K t06-tit-5 sub
\VUsaut{3.6mm}
\VUpk{10.2pt}{40}{غرفة الطعام.}
\VUsaut{3.0mm}

%%K t06-c3-01-1 p
1. --- قُرِع جرس العشاء. واجتمعت الأسرة كلها، ما عدا ماري، حول المائدة.
وغرفة الطعام دافئة دفئًا لطيفًا بفضل \VUgras{مِدفأة} \textsuperscript{(1)} خزفية
ممتازة ذات \VUgras{مِدخنة} \textsuperscript{(2)} رفيعة.

%%K t06-c3-02-1 p
2. --- وليس في هذه الغرفة أثاث كثير. فعلى امتداد الجدار، فوق
\VUgras{المقعد} \textsuperscript{(4)}، تُعلَّق \VUgras{نافورة} \textsuperscript{(3)}
صغيرة للزينة. وقربها \VUgras{الخزانة الجانبية} \textsuperscript{(5)} التي توضع
عليها الأطباق الباردة والحلويات وسائر أدوات المائدة غير اللازمة الآن.

%%K t06-c3-03-1 p
3. --- ففي الرف الأعلى نرى \VUgras{دجاجة مشوية} \textsuperscript{(7)}
و\VUgras{سمكة} \textsuperscript{(8)} و\VUgras{سلطانية} \textsuperscript{(10)} من
\VUgras{الفاكهة} \textsuperscript{(9)}. وتحتها \VUgras{الجبن} \textsuperscript{(11)}
و\VUgras{الخبز} \textsuperscript{(12)} و\VUgras{حامل التوابل} \textsuperscript{(13)}
وفيه كل ما يلزم لتتبيل السلطة: الزيت والخل و\VUgras{الملح} \textsuperscript{(14)}
و\VUgras{الفلفل} \textsuperscript{(15)}. وأخيرًا، في الرف الأسفل،
\VUgras{كعكة} \textsuperscript{(17)} إلى جانب \VUgras{وعاء الخردل}
\textsuperscript{(16)}.

%%K t06-c3-04-1 p
4. --- وساعة غرفة الطعام، وتسمى \VUgras{ساعة الحائط} \textsuperscript{(20)}،
غريبة الشكل جدًا. وهي مركبة في إطار خشبي يحمل أيضًا \VUgras{بارومترًا}
\textsuperscript{(19)} و\VUgras{ترمومترًا} \textsuperscript{(18)}.

%%K t06-tit-6 sub
\VUsaut{4.6mm}
\VUpk{10.2pt}{40}{تقديم العشاء.}
\VUsaut{3.0mm}

%%K t06-c3-05-1 p
5. --- وجدران الغرفة كلها مكسوة بـ\VUgras{ألواح خشبية} \textsuperscript{(21)}
من البلوط المشمّع. و\VUgras{ستارة الباب} \textsuperscript{(22)} القماشية تسدّ
سدًا حسنًا الباب المؤدي إلى الشرفة. وهناك تذهب الأسرة لشرب القهوة بعد
العشاء. و\VUgras{الفناجين} \textsuperscript{(24)} و\VUgras{الصحون}
\textsuperscript{(25)} موضوعة من قبلُ على \VUgras{خزانة الخزف} \textsuperscript{(23)}.

%%K t06-c3-06-1 p
6. --- وفوق المائدة \VUgras{مصباح معلق} \textsuperscript{(26)} مثبّت في السقف؛
ونوره الشديد يملأ غرفة الطعام كلها في المساء.

%%K t06-c3-07-1 p
7. --- ولننظر الآن في \VUgras{مائدة الطعام} \textsuperscript{(27)} نفسها. حين
دخلت الأسرة كانت المائدة مُعَدَّة. فعلى \VUgras{المفرش} \textsuperscript{(28)}
وُضع أمام كل آكل \VUgras{صحن} \textsuperscript{(33)} و\VUgras{منديل}
\textsuperscript{(29)} و\VUgras{ملعقة} \textsuperscript{(34)} و\VUgras{شوكة}
\textsuperscript{(35)} و\VUgras{سكين} \textsuperscript{(36)} و\VUgras{كأس}
\textsuperscript{(32)}. وكانت هناك أيضًا \VUgras{زجاجات} \textsuperscript{(30)} للخمر
و\VUgras{دورق} \textsuperscript{(31)} للماء.

%%K t06-c3-08-1 p
8. --- وقد جاء \VUgras{الخادم} \textsuperscript{(39)} أولًا بـ\VUgras{وعاء
الحساء} \textsuperscript{(37)} الذي لا يزال الحساء يتصاعد منه بخارًا. وهو الآن
يقدّم \VUgras{الطبق} \textsuperscript{(38)} الأول، وهو من اللحم. ثم يدير ما سبق
أن ذكرناه؛ وأخيرًا، بعد \VUgras{الحلوى}، يغتنم انتقال أصحابه إلى الشرفة
ليرفع المائدة ويعيد غرفة الطعام إلى نظامها.

%%K t06-c3-08-2 p
\textit{ملاحظة.} --- \textit{ألفاظ الطعام المتداولة لم تُعطَ هنا إلا في
جملتها. وستُستكمل القائمة في لوحتَي «الفندق» (رقم ١٣) و«السوق»
(رقم ١٥).}

%%K t06-tit-7 sub
\VUsaut{4.6mm}
\VUpk{10.2pt}{40}{المطبخ.}
\VUsaut{3.0mm}

%%K t06-c4-01-1 p
1. --- وفي المطبخ الذي نلقي فيه نظرة من الباب الموارب، نجد فوضى
طريفة. فأمام \VUgras{الموقد} \textsuperscript{(1)} حيث تُطقطق \VUgras{النار}
\textsuperscript{(3)}، نُصب \VUgras{جهاز الشيّ} \textsuperscript{(5)}. و\VUgras{شواء}
\textsuperscript{(7)} جميل على \VUgras{السفّود} \textsuperscript{(6)} يوشك أن ينضج.

%%K t06-c4-02-1 p
2. --- و\VUgras{الكير} \textsuperscript{(8)} و\VUgras{مجرفة الفحم}
\textsuperscript{(9)}، وقد استُعملا لتوّهما، تُركا على الأرض مع
\VUgras{الأحطاب} \textsuperscript{(10)}.

%%K t06-c4-03-1 p
3. --- وأعلى المدفأة مرتب: \VUgras{الفانوس} \textsuperscript{(11)}،
و\VUgras{حامل علبة الثقاب} \textsuperscript{(13)} وفيه \VUgras{أعواد الثقاب}
\textsuperscript{(12)}، و\VUgras{الشمعدان} \textsuperscript{(14)}،
و\VUgras{شمعدان اليد} \textsuperscript{(15)} بـ\VUgras{شمعته} \textsuperscript{(16)}،
كلها في مواضعها. لكن \VUgras{دلو الفحم} \textsuperscript{(18)} الكبير، الملآن
بـ\VUgras{الفحم} \textsuperscript{(17)} الذي نزلت \VUgras{الطاهية}
\textsuperscript{(23)} لتوّها إلى القبو لتجلبه، تُرك في وسط المطبخ.

%%K t06-c4-04-1 p
4. --- والحق أن المسكينة عليها أن تعجل ليكون الغداء جاهزًا في وقته. وهي
الآن عند \VUgras{الموقد} \textsuperscript{(19)} تقلّب \VUgras{صلصة}
\textsuperscript{(24)} لا ينبغي أن تحترق.

%%K t06-c4-05-1 p
5. --- وفي \VUgras{الفرن} \textsuperscript{(20)} تنضج كعكة صنعتها للأطفال
لعصريتهم. وفوق الموقد يُعلَّق \VUgras{المغرفة} \textsuperscript{(21)}
و\VUgras{المصفاة} \textsuperscript{(22)}.

%%K t06-tit-8 sub
\VUsaut{4.0mm}
\VUpk{10.2pt}{40}{أواني المطبخ.}
\VUsaut{3.0mm}

%%K t06-c4-06-1 p
6. --- و\VUgras{طقم القدور} \textsuperscript{(25)} المعلق في آخر الغرفة نظيف
جدًا. فـ\VUgras{الطناجر} \textsuperscript{(26)} النحاسية الجميلة
و\VUgras{إبريق الحليب} \textsuperscript{(27)} و\VUgras{المصفاة} \textsuperscript{(28)}
و\VUgras{المبشرة} \textsuperscript{(29)} قد جُليت بعناية.

%%K t06-c4-07-1 p
7. --- و\VUgras{علبة الملح} \textsuperscript{(30)} على الخزانة الصغيرة إلى جانب
\VUgras{إبريق الشاي} \textsuperscript{(31)} و\VUgras{دَنّ الزيت} \textsuperscript{(32)}.
و\VUgras{الدرج} \textsuperscript{(33)} فيه، على الأرجح، أدوات المائدة وسكاكين
المطبخ.

%%K t06-c4-08-1 p
8. --- و\VUgras{المكنسة} \textsuperscript{(34)} و\VUgras{منفضة الريش}
\textsuperscript{(35)} في ركن. وقد استعملت الطاهية \VUgras{كتلة التقطيع}
\textsuperscript{(36)} لتوّها؛ وتركتها قرب النافذة.

%%K t06-c4-09-1 p
9. --- و\VUgras{منشفة يد} \textsuperscript{(37)} معلقة على الجدار إلى جانب
\VUgras{القمع} \textsuperscript{(38)}. وفي هذا الجانب من المطبخ يُحفظ
\VUgras{الشوّاية} \textsuperscript{(39)} و\VUgras{الساطور} \textsuperscript{(40)}
و\VUgras{لوح التقطيع} \textsuperscript{(41)}. ثم، فوق الساطور، على
\VUgras{رفّ} \textsuperscript{(42)}، \VUgras{أوانٍ} \textsuperscript{(43)}،
و\VUgras{قِدر الطبخ} \textsuperscript{(44)} بـ\VUgras{غطائها} \textsuperscript{(45)}،
و\VUgras{المرجل} \textsuperscript{(46)}. و\VUgras{المقلاة} \textsuperscript{(47)}
معلقة على مقربة.

%%K t06-c4-10-1 p
10. --- ولا يلمع في هذا الركن إلا \VUgras{الصنبور} \textsuperscript{(48)}
النحاسي. و\VUgras{المجلى} \textsuperscript{(49)}، الذي يقوم عليه \VUgras{الجرّة}
\textsuperscript{(50)} الكبيرة، لا بد أنه مريح جدًا بخزانته الصغيرة، حيث
يُحفَظ \VUgras{برميل الخل} \textsuperscript{(51)} بـ\VUgras{صنبوره}
\textsuperscript{(52)}، و\VUgras{صندوق النفايات} \textsuperscript{(53)}،
و\VUgras{الصحون الوسخة} \textsuperscript{(54)}.

%%K t06-tit-9 sub
\VUsaut{4.0mm}
\VUpk{10.2pt}{40}{ما يُحفظ في المطبخ أيضًا.}
\VUsaut{3.0mm}

%%K t06-c4-11-1 p
11. --- و\VUgras{الجفنة} \textsuperscript{(55)} الكبيرة التي تُغسل فيها
\VUgras{الخضر} \textsuperscript{(58)}، و\VUgras{المرجل} \textsuperscript{(56)}
الحديدي القائم على \VUgras{الأرضية المبلَّطة} \textsuperscript{(57)}، من تلك
الخزانة كذلك على الأرجح.

%%K t06-c4-12-1 p
12. --- ولا تعوز الطاهية المواقد بالتأكيد، فها هنا على طرف المائدة
\VUgras{موقد غاز} \textsuperscript{(59)} عليه ماء يسخن في طنجرة صغيرة.
و\VUgras{البخار} \textsuperscript{(60)} المتصاعد منها ينبئنا أنه لا بد يغلي.
وهذا الماء للقهوة على الأرجح، فها هنا في الطرف الآخر من المائدة
\VUgras{ركوة القهوة} \textsuperscript{(64)} و\VUgras{مطحنة القهوة}
\textsuperscript{(63)}.

%%K t06-c4-13-1 p
13. --- وموقد الغاز موصول بـ\VUgras{فوهة الغاز} \textsuperscript{(62)}
بـ\VUgras{أنبوب مطاطي} \textsuperscript{(61)} طويل.

%%K t06-c4-14-1 p
14. --- و\VUgras{العاملة بالمياومة} \textsuperscript{(66)}، التي تأتي لتعين
الطاهية، تهيّئ خضر العشاء. فإذا قُشّرت وغُسلت وضعتها في \VUgras{القدر}
\textsuperscript{(65)} إلى أن يحين وقت طبخها.

%%K t06-tit-10 sub
\VUsaut{4.0mm}
{\centering\fontsize{10.2pt}{10.2pt}\selectfont\textls[40]{\scshape الحمّام.} \textit{(غرفة الاستحمام.)}\par}
\VUsaut{3.0mm}

%%K t06-c5-01-1 p
1. --- ولم يبق إلا تطفّل واحد لنعرف غرف البيت الرئيسية: أن ننظر في
تجهيزات الحمّام. ولن يطول ذلك، فهو غير واسع، وسنرى سريعًا أن
\VUgras{حوض الاستحمام} \textsuperscript{(1)} له \VUgras{سخّان} \textsuperscript{(2)}
جيد، وأن \VUgras{صحن الصابون} \textsuperscript{(3)} في متناول يد المستحم، وأن
\VUgras{حامل المناشف} \textsuperscript{(5)} لا بد أنه يسهّل تجفيف
\VUgras{المناشف} \textsuperscript{(4)}.

%%K t06-c5-02-1 p
2. --- وهذه الغرفة جدرانها مكسوة بـ\VUgras{بلاط مزجَّج} \textsuperscript{(6)}،
وهو ما أتاح تركيب \VUgras{دُشّ} \textsuperscript{(7)} دون خشية أن يفسدها الماء
المتناثر أثناء الاستعمال. و\VUgras{طست} \textsuperscript{(8)} زنكي صغير لحمّام
القدمين يتمّ الأثاث.
\VUsaut{6.0mm}
\VUfilet{22mm}
